\chapter{Results and Discussion}
\label{chap:evaluation}

This chapter reports and interprets the empirical results of the trained
artifacts: predictive performance against non-learned baselines, calibration
behaviour, feature attribution and error analysis, and an ablation of the
human-actor filter.

The central result is stated up front, because it shapes how the rest of the
chapter should be read. On the held-out slice as a whole, a single observable
signal --- the number of days since the last human commit --- already ranks
inactivity almost as well as a sixty-feature gradient-boosted model. The learned
models are nevertheless worth their complexity, but not for the reason an
aggregate metric would suggest: their advantage concentrates on the repositories
that are still active at the observation time, which is precisely the population
for which an early-warning signal has any decision value at all.

\section{Predictive Performance}
\label{sec:results-performance}

Table~\ref{tab:model-comparison} reports held-out test metrics for the three
learned model families in both feature regimes, together with two non-learned
baselines. All learned models were trained on the same time-aware split
(\(75/15/10\); \(n_\text{test}=5{,}000\) labelled observations, inactive base
rate \(0.661\)) and calibrated on the validation slice with the same histogram
calibrator, so the comparison isolates the effect of model family rather than
data or calibration differences. The combined quality score summarizes ranking
skill (ROC-AUC) and calibration (Brier) into a single comparison figure and is
not intended as a standalone proof. All reported metrics are defined formally in
Section~\ref{sec:method-metric-definitions}.

\begin{table}[ht]
  \centering
  \caption{Held-out test metrics by model family and feature regime, with the two
  trivial baselines. Higher is better for ROC-AUC, F1, and quality; lower is
  better for Brier and ECE. The best value within each regime is shown in bold.
  The recency baseline is defined on the \(4{,}556\) test rows that have an
  observable commit history; see the text for the like-for-like comparison.}
  \label{tab:model-comparison}
  \small
  \setlength{\tabcolsep}{3.5pt}
  \begin{tabular}{lrrrrr}
    \toprule
    Model & ROC-AUC & Brier & ECE & F1 & Quality \\
    \midrule
    \texttt{logistic-regression-full-history} & 0.949 & 0.088 & 0.047 & 0.911 & 0.782 \\
    \texttt{xgboost-full-history}             & \textbf{0.958} & \textbf{0.078} & 0.035 & \textbf{0.916} & \textbf{0.810} \\
    \texttt{neural-net-full-history}          & 0.956 & 0.081 & \textbf{0.028} & 0.913 & 0.803 \\
    \midrule
    \texttt{logistic-regression-cold-start}   & 0.932 & 0.100 & 0.044 & 0.897 & 0.740 \\
    \texttt{xgboost-cold-start}               & \textbf{0.953} & \textbf{0.082} & 0.033 & \textbf{0.909} & \textbf{0.798} \\
    \texttt{neural-net-cold-start}            & 0.951 & 0.084 & \textbf{0.027} & 0.908 & 0.791 \\
    \midrule
    \texttt{recency-baseline} (days since commit)   & 0.957\rlap{$^{*}$} & --- & --- & 0.918 & --- \\
    \texttt{commit-volume-baseline} (\(-\)commits\_365d) & 0.897 & --- & --- & 0.904 & --- \\
    \bottomrule
  \end{tabular}
\end{table}

Among the learned models the ordering is stable across both regimes. The
gradient-boosted trees (\texttt{xgboost\sdash full\sdash history}) achieve the strongest
ranking and probabilistic accuracy (ROC-AUC \(0.958\), Brier \(0.078\)), ahead of
the multilayer perceptron (\(0.956\) / \(0.081\)) and logistic regression
(\(0.949\) / \(0.088\)). The neural network was included as a deliberate test of
whether a higher-capacity non-linear model is warranted for this tabular feature
space. It does not improve on the gradient-boosted trees on ranking or Brier
score in either regime, while adding training cost and weaker interpretability;
its one advantage is the lowest expected calibration error. This is consistent
with the broader finding that tree-based models remain a strong default on
moderate-sized tabular problems~\cite{grinsztajn2022trees}. Accordingly, logistic
regression is retained for its transparent coefficient attribution and XGBoost as
the deployed runtime scorer, while the neural network is reported as a
tested-and-rejected alternative rather than promoted into the scoring ensemble.

The cold-start regime is only marginally weaker than the full-history regime for
the non-linear models: \texttt{xgboost\sdash cold\sdash start} reaches ROC-AUC \(0.953\)
against \(0.958\) for its full-history counterpart, despite using 17 rather than
60 features and, in particular, despite lacking the commit-window features that
dominate the full-history model's gain (Section~\ref{sec:results-attribution}).
Only the linear model degrades appreciably (\(0.932\) against \(0.949\)). The
regimes are therefore not as far apart as the feature counts suggest, because the
cold-start vector retains \texttt{last\_push\_age\_days}, a recency signal that is
nearly as informative as the reconstructed commit history. This appears to qualify
the expectation stated in Section~\ref{sec:method-prediction-regimes} that the
full-history model would be clearly superior: for the deployed tree model it is
better, but only by \(0.005\) ROC-AUC. That reading does not survive a closer
look. Section~\ref{sec:results-history-ablation} shows the aggregate gap to be an
artifact of the slice composition, and the regimes to separate by \(0.030\)
ROC-AUC once the already-quiet majority is removed.

\subsection{Comparison Against Non-Learned Baselines}
\label{sec:results-baselines}

Because the inactivity label is itself activity-derived, a learned model is only
worth reporting if it improves on a trivial rule built from a single observable
signal (Section~\ref{sec:method-model-training}). Two such rules are evaluated on
the same held-out slice: ranking by days since the last human commit, and ranking
by negated annual commit volume.

The result is sobering, and it is the most important quantitative finding of this
thesis. The recency rule attains ROC-AUC \(0.957\) and an F1 score of \(0.918\),
which is not merely competitive with the deployed \texttt{xgboost\sdash full\sdash history}
model but nominally exceeds its F1 of \(0.916\). The comparison is not quite
like-for-like: \texttt{days\_since\_last\_commit} is undefined for the \(444\)
test rows whose repositories have no observable record before the observation date
(Section~\ref{sec:method-undefined-features}), so the baseline is scored on
\(4{,}556\) rows while the model is scored on all \(5{,}000\). Re-scoring the
deployed model on exactly those \(4{,}556\) rows raises it to ROC-AUC \(0.974\)
against the baseline's \(0.957\). The model therefore does outrank the trivial
rule, but by \(0.017\) ROC-AUC rather than by the wide margin that an aggregate
comparison against a weaker baseline would suggest. The commit-volume rule is
clearly worse in aggregate (\(0.897\)).

Two conclusions follow. First, the aggregate held-out metric is a poor instrument
for judging this artifact: most of the achievable ranking skill on a slice whose
inactive base rate is \(0.661\) is already contained in one recency signal, and a
sixty-feature model recovers only a modest increment on top of it. Reporting the
headline ROC-AUC alone would overstate what the learned model contributes.
Second, this is exactly the situation the active-at-\(t\) evaluation was designed
to expose (Section~\ref{sec:method-evaluation-strategy}). Among repositories that
are already quiet at \(t\), predicting continued inactivity is close to trivial,
and both the model and the recency rule do it well; the aggregate metric is
dominated by that easy majority. The decision-relevant question is what happens on
the repositories that are still active, and there the picture separates
sharply, as Section~\ref{sec:results-slices} shows.

\subsection{How Much Is Reconstructed History Worth?}
\label{sec:results-history-ablation}

The same instrument answers the feature-availability question posed in
Section~\ref{sec:method-prediction-regimes}. Comparing the two gradient-boosted
models isolates the contribution of the 43 reconstructed historical features,
since the cold-start model differs from the full-history model only by their
absence. Table~\ref{tab:history-ablation} reports every arm on identical rows:
where the recency signal is undefined the row is dropped from \emph{all} arms
rather than imputed for the models alone.

\begin{table}[htbp]
  \centering
  \caption{Held-out ROC-AUC by feature-set size, on identical row sets.}
  \label{tab:history-ablation}
  \begin{tabular}{lrrrr}
    \toprule
    Subset & \(n\) & Full-history (60) & Cold-start (17) & Best trivial rule (1) \\
    \midrule
    Recency-defined & 4{,}556 & 0.974 & 0.970 & 0.958 \\
    active@\(t\)    & 1{,}794 & 0.905 & 0.875 & 0.844 \\
    \bottomrule
  \end{tabular}
\end{table}

In aggregate, reconstructed history is worth \(0.005\) ROC-AUC, and the 17
snapshot features are worth \(0.013\) over the single best trivial signal. Read
alone, those numbers say the GH Archive pipeline was not worth building. On the
active-at-\(t\) subset the comparison inverts: history is worth \(0.030\), six
times its aggregate value, and the snapshot features \(0.083\). The ordering of
the three feature sets --- one signal, seventeen, sixty --- is strictly monotone
where early warning matters and almost flat where it does not.

This is the third independent appearance of one property of the data. The learned
model beats the trivial rule by \(0.017\) in aggregate against \(0.061\) on
active-at-\(t\); reconstructed history by \(0.005\) against \(0.030\); the
snapshot feature set by \(0.013\) against \(0.083\). Each aggregate comparison is
compressed by the same cause: a slice in which two thirds of repositories are
already quiet rewards a predictor for restating the present, so any additional
information --- more features, more history, a learned decision boundary --- is
visible only once that easy majority is removed. The aggregate ROC-AUC of this
dataset measures activity persistence more than predictive skill, and no design
decision in this thesis should be judged against it.

\section{Calibration and Threshold Selection}
\label{sec:results-calibration}

Because the score is intended to be read as a probability rather than only as a
ranking, calibration is evaluated alongside discrimination. Every model's raw
outputs are mapped through a histogram calibrator fitted on the validation slice
and never on the test slice, so the reported calibration reflects held-out
behaviour. For the deployed \texttt{xgboost\sdash full\sdash history} model this yields a
Brier score of \(0.078\), a log loss of \(0.246\), and an expected calibration
error of \(0.035\) at the selected decision threshold of \(0.51\), where the
operating point attains precision \(0.90\) and recall \(0.93\). Among the
full-history models the neural network attains the lowest expected calibration
error (\(0.028\)) and logistic regression the highest (\(0.047\)), so the ranking
by calibration does not follow the ranking by discrimination.

The reliability of the calibrated probabilities is well behaved: the empirical
inactivity rate rises monotonically across the ten calibration bins, from
\(0.032\) in the lowest predicted-probability band to \(0.988\) in the highest, so
higher scores do correspond to higher observed inactivity frequencies rather than
only to a higher rank. Two caveats apply. First, the two extreme bands hold
\(65.9\%\) of the validation mass, so mid-range probabilities are supported by
comparatively few samples and are correspondingly less certain; the per-prediction
confidence surfaced by the demonstrator (Chapter~\ref{chap:demonstrator}) reflects
exactly this bin-level support. Second, calibration is conditional on the sampled
prevalence rather than a natural deployment prevalence
(Section~\ref{sec:method-limitations}), so the absolute probabilities would
require recalibration before being read as live population probabilities.

\section{Feature Importance and Error Analysis}
\label{sec:results-features}

\subsection{Feature Attribution}
\label{sec:results-attribution}

Feature attribution is read from the logistic-regression model, which is retained
precisely for its transparent coefficients (Section~\ref{sec:results-performance}).
Because the inputs are standardized, the coefficient magnitudes are directly
comparable, and their signs indicate the direction of the association with
inactivity. The strongest inactivity-\emph{lowering} signal is the breadth of
sustained development (\texttt{active\_commit\_months\_365d}, standardized
coefficient \(-0.79\)), followed by recent maintainer integration
(\texttt{has\_recent\_pr\_merge\_flag}, \(-0.39\)) and the presence of any issue
activity at all (\texttt{has\_issue\_activity\_365d}, \(-0.26\)). The strongest
inactivity-\emph{raising} signals are the two staleness measures, days since the
last push and days since the last commit (each \(+0.56\)). These directions match
the intended semantics of the features (Appendix~\ref{appendix:feature-reference}).

Two coefficients resist that reading. A larger open-issue backlog is associated
with \emph{lower} inactivity (\texttt{open\_issues\_log1p}, \(-0.41\)), and a
larger star count with \emph{higher} inactivity (\texttt{stars\_log1p},
\(+0.25\)). Both are plausibly confounded rather than causal: an open backlog
presupposes a user base that files issues, and popularity is correlated with
project age within the sampled frame. These coefficients are therefore used to
explain \emph{direction} for a given repository, while ranking performance is
carried by the gradient-boosted model.

These attributions are only interpretable because of the encoding correction of
Section~\ref{sec:method-undefined-features}. While an undefined response latency was
stored as zero, the responsiveness medians ranked among the strongest inactivity
signals, inviting the reading that slow-responding projects fail; in fact a zero
latency marked projects with \emph{no} issues or pull requests at all. After the
correction those medians no longer appear among the leading coefficients, and the
absence of activity is carried explicitly by \texttt{has\_issue\_activity\_365d} and
\texttt{has\_pr\_activity\_365d}.

The gradient-boosted model exposes a complementary, non-linear view through its
gain-based feature importances (normalized to sum to one over the \(42\) of \(60\)
features that receive at least one split). Its ranking is dominated by recent
commit volume: \texttt{commits\_365d} (\(0.42\)) and \texttt{commits\_90d}
(\(0.20\)) together account for over three-fifths of the total gain, followed by
\texttt{active\_commit\_months\_365d} (\(0.08\)) and
\texttt{has\_recent\_pr\_merge\_flag} (\(0.07\)); \texttt{last\_push\_age\_days},
\texttt{opened\_prs\_90d}, \texttt{days\_since\_last\_commit}, and
\texttt{merged\_prs\_90d} contribute the next tier. The two views agree on the
substantive story --- sustained recent development activity is the dominant signal
--- while differing in emphasis: the linear model leans on the breadth of active
commit months, whereas the gradient-boosted model concentrates gain on the raw
recent-commit-volume windows. That concentration also explains the aggregate near-parity of
the cold-start model observed in Section~\ref{sec:results-performance}: the
cold-start vector omits both commit windows, yet \texttt{last\_push\_age\_days}
substitutes for them at little cost \emph{on a slice where most repositories are
already quiet}. Where they are not, the substitution is a poorer one
(Section~\ref{sec:results-history-ablation}).

\subsection{Robustness to Repository Recurrence}
\label{sec:results-robustness}

Under the repository-disjoint split (Section~\ref{sec:impl-grouped-split}), where
every snapshot of a repository is confined to one partition, held-out ROC-AUC
decreases from \(0.958\) to \(0.950\) and the Brier score rises from \(0.078\) to
\(0.079\). The small gap bounds the optimism attributable to a repository
appearing in both training and test slices and indicates that the discrimination
generalizes to unseen repositories rather than resting on repository
memorization. This check matters more than it did previously: because the labelled
dataset now spans ten equally sized quarterly cohorts, the chronological test
slice consists entirely of the final observation date, so every test repository
also appears in training at an earlier date.

\subsection{Slice Evaluation}
\label{sec:results-slices}

A single aggregate metric can average a strong subgroup with a weak one. The
held-out predictions of the deployed \texttt{xgboost\sdash full\sdash history} model are
therefore re-scored within subgroups that vary in ways relevant to deployment; the
full per-subgroup table is given in Appendix~\ref{appendix:eval-tables}
(Table~\ref{tab:slice-metrics}). Three findings follow.

First, and most importantly for the activity-persistence threat
(Section~\ref{sec:validity-internal}), the model separates outcomes well
\emph{among repositories that are active at the observation time}: on the
active-at-\(t\) subset (\(n=1{,}794\), inactive rate only \(0.225\)) it attains
ROC-AUC \(0.905\). On this same subset the recency baseline reaches only
\(0.792\) and the commit-volume baseline \(0.844\). The learned model therefore
holds a margin of \(0.061\) over the best trivial rule exactly where early-warning
skill matters, against a margin of \(0.017\) in aggregate
(Section~\ref{sec:results-baselines}). This is the central justification for the
artifact: its value is not that it ranks a mixed population better than a recency
rule, but that it anticipates decline in repositories that still look healthy.
The operating point on this subset is nevertheless conservative --- recall
\(0.49\) at precision \(0.74\) --- so the model flags fewer than half of the
repositories that will fall quiet, and the score is better used as a ranking for
triage than as an automatic decision.

Second, the aggregate ROC-AUC masks a much weaker low-popularity subgroup: those
repositories are ranked at only \(0.846\), against \(0.972\) for high-popularity
and \(0.964\) for medium-popularity ones. Scores for low-visibility projects
deserve more caution and are exactly where the per-prediction confidence should be
low. Note that ``low popularity'' here is relative to a seed frame that already
requires at least 100 stars, so these are not obscure repositories in absolute
terms.

Third, full history helps less than the two-regime design assumed. Repositories
with reconstructable history are ranked at \(0.921\), against \(0.929\) for
cold-start-like repositories --- a reversal of the expected ordering. The
cold-start-like subgroup has a far higher inactive rate (\(0.924\) against
\(0.350\)), so its ranking problem is easier; the comparison across subgroups is
therefore confounded by prevalence and should not be read as evidence that
reconstructed history is unhelpful. It does, however, reinforce the finding of
Section~\ref{sec:results-performance} that the two regimes are closer in skill
than their feature counts suggest.

\subsection{Error Analysis}
\label{sec:results-error-analysis}

The comparison against explicit trivial baselines is now part of the reported
results (Section~\ref{sec:results-baselines} and the slice analysis above); the
systematic formative case review that motivated the feature-set revision
(Section~\ref{sec:method-iterative-development}) remains a complementary
qualitative analysis and is noted as a remaining item in
Chapter~\ref{chap:conclusion}. The quantitative results identify two residual
weaknesses. The dominant one is low-visibility repositories with sparse public
signals, which the confidence layer is designed to flag rather than hide. The
second is recall on the active-at-\(t\) subset: the model ranks impending decline
well but, at the F1-selected threshold, commits to flagging only about half of it.
Both weaknesses argue for using the score as a prioritization signal rather than as
a gate.

\section{Ablation: The Human-Actor Filter}
\label{sec:results-bot-ablation}

Commit- and contributor-based features and the maintenance label are computed from
human actors only; automated (bot) accounts are excluded
(Section~\ref{sec:method-feature-engineering}). Because this is a design choice
that could plausibly be omitted, it is tested directly by an ablation: the full
5{,}000-repository dataset was rebuilt with the human-actor filter disabled through
a single build flag, with every other setting identical, yielding 50{,}000
snapshots that align one-to-one with the filtered build. Three effects are
examined.

\textbf{Label corruption.} Disabling the filter flips 357 of 50{,}000 labelled
snapshots (\(0.71\%\)) from inactive to active, and \emph{none} in the opposite
direction. The effect is strictly one-directional: a bot commit or merge inside the
future window can satisfy the two-of-four maintenance rule, so an otherwise dormant
repository is relabelled as maintained, whereas bots never make an active
repository look inactive. Omitting the filter would therefore bias the target by
systematically \emph{masking} inactivity, which is the error direction that matters
most for a risk signal.

\textbf{Feature inflation.} On the observation-time side, bots inflate the activity
features substantially: with the filter disabled, \texttt{commits\_365d} changes on
7{,}858 of 50{,}000 snapshots (\(15.7\%\)) by a mean of \(+938\) commits, though
the median change is only \(+80\), so the mean is driven by a small number of
heavily automated repositories. \texttt{commits\_90d} rises by a mean of \(+388\)
and \texttt{contributors\_365d} by \(+1.4\) distinct contributors, while
pull-request \emph{count} features are unaffected because they are not
actor-filtered. A small number of bot-driven repositories would thus dominate the
activity signal.

\textbf{Predictive effect.} A naive comparison, in which each model is scored
against the labels of the dataset it was trained on, would show bots to be nearly
harmless --- but it is misleading, because the unfiltered model is then scored
against a partially corrupted target. A cleaner test trains on the bot-contaminated
\emph{features} while evaluating against the clean, human-filtered \emph{labels}
(Appendix~\ref{appendix:eval-tables}, Table~\ref{tab:bot-ablation-crosstest}): the
contaminated features are consistently worse across all four families, though by a
small margin (ROC-AUC falls by \(0.0006\)--\(0.0011\); Brier rises by up to
\(0.0010\)).

\textbf{Interpretation.} The filter's justification is construct validity rather
than discrimination. Its predictive cost when omitted is small, since the
gradient-boosted models are largely robust to the minority of bot-inflated
repositories, but it is consistently in the expected direction, and, more
importantly, the filter removes a one-directional label bias that a headline metric
computed on the corrupted data does not reveal. The human-actor filter is therefore
retained.
